\appendix
\chapter{Prompts and Reproduction Reference}
\label{chap:appendix}

\section{Complete Prompt Reference}

For convenience, this section collects every prompt referenced in the
body of the report in one place. All are transcribed verbatim (module
docstrings/comments removed) from
\url{https://github.com/AnnaRemi/ENSIMAG-AI-M2-final-project/tree/agent/consolidate-four-method-benchmarks}.

\begin{lstlisting}[style=prompt, caption={SUQL baseline \answerfn{} system prompt (\texttt{project SUQL/baseline/suql\_engine.py}).}]
You are a recall-first semantic filter for movie reviews.

Decide whether the review provides evidence that the movie
satisfies the question.

Recall is the primary objective, but do not accept every
review:
- Answer YES when there is any concrete review-specific
  evidence: direct wording, a synonym/paraphrase, a described
  example, or a reasonable implication.
- Give borderline evidence the benefit of the doubt when it
  specifically bears on the condition. The evidence need not
  use the question's exact words.
- This is evidence retrieval, not sentiment scoring: an
  explicit mention still counts when it is negated, qualified,
  quoted, or used critically.
- Answer NO when there is no review-specific support, the
  connection is only the movie's genre/topic, the text is
  generic praise or criticism, or the review never discusses
  the condition itself.
- Do not infer evidence merely because the condition is not
  contradicted.

Return exactly YES or NO. Do not include reasoning,
punctuation, or extra text.
\end{lstlisting}

\begin{lstlisting}[style=prompt, caption={SUQL semantic-parser system prompt (\texttt{project SUQL/baseline/suql\_engine.py}, \texttt{PARSER\_SYSTEM}), which compiles the NL question into a SUQL query.}]
You are a SUQL (Structured and Unstructured Query Language)
semantic parser.
SUQL extends SQL with two free-text functions:
  answer(free_text_column, 'question') -> 'Yes'|'No'|short string
  summary(free_text_column)             -> a short prose summary

{TABLE_SCHEMA}

Rules:
1. Use plain SQL WHERE clauses for ALL structured predicates
   (year, runtime, director, genres, title, movie_id).
2. Use answer() ONLY for predicates that require understanding
   the review text.
3. Always apply structured filters BEFORE answer() in the
   WHERE clause (efficiency).
4. The SELECT list must include: movie_id, title, year,
   runtime, director, genres. Add summary(review) AS
   review_summary whenever you use answer() on the review.
5. Use LIKE '%value%' for partial text matches on genres and
   director.
6. For "top N" questions without a numeric ranking field, use
   LIMIT N.
7. Output ONLY the raw SQL query -- no markdown fences, no
   explanation.

{FEW_SHOT_EXAMPLES}
\end{lstlisting}

\begin{lstlisting}[style=prompt, caption={Trummer baseline block-join prompt (\texttt{project Trummer/baseline/trummer\_join/operators.py}, \texttt{create\_prompt}).}]
Find every movie in Collection 1 that has a review in
Collection 2 satisfying this predicate:
{predicate}
Act as a recall-first semantic filter, while requiring
review-specific evidence.
- Answer YES for direct wording, synonyms/paraphrases,
  described examples, or reasonable implications.
- Give borderline evidence the benefit of the doubt when it
  specifically bears on the predicate.
- Count explicit mentions even when negated, qualified,
  quoted, or critical; retrieve evidence rather than score
  sentiment.
- Answer NO when support is absent, generic, based only on
  genre/topic, or never discusses the predicate itself.
- Lack of contradiction alone is not evidence.
Collection 1 contains structured movie rows. Collection 2
contains review chunks.
A review belongs to a movie only when tconst is exactly
equal to movie_id.
Return one decision per line as: movie_id: YES or
movie_id: NO.
Do not include confidence, evidence, explanations, or prose.
Copy each movie_id exactly from Collection 1.
Do not add prose, JSON, markdown, or a completion marker.

Collection 1:
{block_1_rows}

Collection 2:
{block_2_rows}

Decisions:
\end{lstlisting}

\begin{lstlisting}[style=prompt, caption={SUQL~V1 cheap Stage~1 single-token scoring prompt (\texttt{project SUQL/v1/scorer.py}).}]
Act as a high-recall first-pass semantic filter. Decide
whether this review contains review-specific evidence for a
Yes answer. Count direct wording, synonyms, described
examples, and reasonable implications as evidence. Give
borderline but condition-specific evidence the benefit of
the doubt. Do not discard an explicit mention because it is
negated, qualified, quoted, or critical: this is evidence
retrieval, not sentiment scoring. Do not count genre/topic
alone, generic praise or criticism, or mere lack of
contradiction as evidence.
Return exactly one token: 1 for Yes, 0 for No.

Question: {question}
Review: {review[:1800]}

Answer:
\end{lstlisting}

\begin{lstlisting}[style=prompt, caption={SUQL~V1 expensive Stage~2 verification prompt (\texttt{project SUQL/v1/cascade\_filter.py}, \texttt{ANSWER\_SYSTEM}).}]
You are the final recall-first semantic filter for movie
reviews.

Decide whether the review provides evidence that the movie
satisfies the question.

Recall is the primary objective, but do not accept every
review:
- Answer YES for any concrete review-specific support,
  including direct wording, synonyms, described examples, and
  reasonable implications.
- Give borderline evidence the benefit of the doubt when it
  specifically bears on the condition; exact wording is
  unnecessary.
- This is evidence retrieval, not sentiment scoring: an
  explicit mention still counts when it is negated, qualified,
  quoted, or used critically.
- Answer NO when support is absent, merely generic, based only
  on genre/topic, or never discusses the condition itself.
  Lack of contradiction alone is not evidence.

Return exactly YES or NO. Do not include reasoning,
punctuation, or extra text.
\end{lstlisting}

\begin{lstlisting}[style=prompt, caption={Trummer~V1 cheap Stage~1 batched prompt (\texttt{project Trummer/v1/trummer\_join/cascade.py}, \texttt{\_cheap\_batch\_prompt}).}]
Classify every explicit candidate pair below.
Predicate: {predicate}
{CHEAP_EVIDENCE_INSTRUCTIONS}
Return one line per candidate as CANDIDATE_ID: YES, NO, or
UNCERTAIN.
Do not include confidence, evidence, explanations, or prose.
Return every candidate exactly once and do not explain.

CANDIDATE_{n}:
Movie: {movie_text}
Review: {review_text}
...

Decisions:
\end{lstlisting}

\begin{lstlisting}[style=prompt, caption={Trummer~V1 expensive Stage~2 batched verification prompt (\texttt{cascade.py}, \texttt{\_expensive\_batch\_prompt}).}]
Evaluate the explicit candidate pairs below.
Predicate: {predicate}
{EXPENSIVE_EVIDENCE_INSTRUCTIONS}
Return one line per candidate as PAIR_ID: YES or PAIR_ID: NO.
Do not include confidence, evidence, explanations, or prose.
Return every candidate exactly once. Do not add prose, JSON,
or markdown.

PAIR_{n}:
Movie: {movie_text}
Review: {review_text}
...

Decisions:
\end{lstlisting}

\section{Calibration Reference Values}
\label{sec:appendix-calibration}

Table~\ref{tab:appendix-thresholds} lists the calibrated cheap-stage
accept/reject log-odds thresholds actually stored in
\texttt{project SUQL/v1/thresholds.json} for a representative sample of
predicates from earlier (pre-\texttt{10q}) calibration runs; the
\texttt{10q} benchmark of Chapter~\ref{chap:experiments} recalibrates
fresh per question rather than reusing this file, but the values
illustrate the typical magnitude of the calibrated band.

\begin{table}[h]
\centering\small
\begin{tabular}{p{9.5cm}rr}
\toprule
Predicate (abbreviated) & Accept $\tau_+$ & Reject $\tau_-$ \\
\midrule
Funny, humorous, witty, or entertaining? & $1.0$ & $-1.0$ \\
Praise acting/performances/cast/characters? & $1.0$ & $-1.0$ \\
Suspense, tension, scares, fear, creepy atmosphere? & $1.0$ & $-1.0$ \\
Praise chemistry, romance, love story, relationships? & $1.0$ & $-1.0$ \\
Exciting, thrilling, intense, action-packed? & $1.0$ & $-1.0$ \\
\bottomrule
\end{tabular}
\caption{Sample calibrated thresholds, \texttt{project SUQL/v1/thresholds.json}.
The symmetric $\pm 1.0$ band corresponds to routing everything with
cheap-stage probability outside roughly $[0.27, 0.73]$ directly, and
escalating the rest.}
\label{tab:appendix-thresholds}
\end{table}

\section{Reproduction: Running the Benchmarks Locally}
\label{sec:appendix-cli}

The repository's runner supports all four methods and all four suite
sizes out of the box; the commands below are the actual entry points
(not a hypothetical interface), taken from the repository's top-level
and \texttt{benchmarks/} README files.

\begin{lstlisting}[style=prompt, caption={Cloning and running the smoke-test suite locally against Ollama.}]
git clone git@github.com:AnnaRemi/ENSIMAG-AI-M2-final-project.git
cd ENSIMAG-AI-M2-final-project
bash benchmarks/1q/run_local.sh --pull-models
\end{lstlisting}

\begin{lstlisting}[style=prompt, caption={Running the full 10-question, 10-repetition suite on the lab's Aker cluster (the exact configuration used to produce Chapter~\ref{chap:experiments}'s tables).}]
bash benchmarks/run_aker.sh \
  --suite 10q \
  --repetitions 10 \
  --methods "suql_baseline suql_v1 trummer_baseline trummer_v1" \
  --cheap-model gemma4:e2b \
  --expensive-model gemma4:26b \
  --pull-models
\end{lstlisting}

Both entry points create a \texttt{.venv}, install the shared Python
dependencies, pull the configured models, verify the local Ollama API,
and run all four implementations end to end, writing
\texttt{aggregate.csv}, \texttt{comparison.csv}, and the four aggregate
plots (quality, time, call-count, and a speedup-vs-call-reduction
scatter, plus per-question versions of the same four plots) to
\texttt{benchmarks/<suite>/outputs/<run-name>/}. Raw model responses and
per-repetition intermediate artifacts are deleted after successful
evaluation by default (pass \texttt{--keep-run-artifacts} to
\texttt{shared/scripts/run\_all.py} directly to retain them for
debugging).
